% part: intuitionistic-logic
% chapter: propositions-as-types
% section: proof-terms

\documentclass[../../../include/open-logic-section]{subfiles}

\begin{document}

\olfileid{int}{pty}{ter}

\olsection{প্রমাণপদ}

\Log{N1i}-তে অনুমানগুলিকে অবমুক্তি-চিহ্ন হিসেবে
চলরাশি দিয়ে (যেমন, $!A^x$) চিহ্নিত করি।
\Log{N2i}-তে প্রতিটি সিকোয়েন্ট
$\Gamma \Sequent !A$-র বাম প্রসঙ্গ~$\Gamma$-তে
চলরাশি--!!{formula} যুগলের তালিকা থাকে। তাই
\Log{N2i}-তে !!a{proof}-এর প্রতিটি সিকোয়েন্ট
এ পর্যন্ত !!{proof} কী !!{prove}d করেছে
(অর্থাৎ~$!A$), সেই তথ্যের সঙ্গে প্রতিটি ধাপে
কোন অনুমানগুলি খোলা আছে (অর্থাৎ~$\Gamma$),
সেই তথ্যও বহন করে। প্রত্যেক সিকোয়েন্টে
$!A$ \emph{কীভাবে} !!{prove}d হয়েছে, সেই
তথ্যও যোগ করলে কেমন হয়? এ জন্য
\emph{পদ-লেবেলযুক্ত} \Log{tN2i} ব্যবস্থা
সংজ্ঞায়িত করি, যেখানে সিকোয়েন্টের রূপ
$\Gamma \Sequent N : !A$। $\Gamma$ ও
$!A$ আগের মতোই: $x: !B$ রূপের যুগলের
একটি সেট, এবং সিকোয়েন্টে শেষ হওয়া
উপ-!!{proof} যে !!{formula}কে~$\Gamma$-র
অনুমানগুলি থেকে অনুসৃত দেখিয়েছে।
$N$ হবে সেই সিকোয়েন্টে শেষ হওয়া
!!{proof}-কে সংকেতায়িত একটি নির্দিষ্ট পদ।

প্রতিটি সিকোয়েন্টে দেওয়া পদগুলি চলরাশি
$x$, $y$, \dots{} থেকে তৈরি করি; এ পর্যন্ত
কোন অনুমিতিগুলি হয়েছে, তা সংকেতায়িত করে
এমন অপেক্ষকচিহ্ন ব্যবহার করি। প্রতিটি
অনুমানবিধির জন্য পৃথক অপেক্ষকচিহ্ন চাই।
!!{introduction} বিধির সংশ্লিষ্ট অপেক্ষকচিহ্নগুলি
“নির্মাতা”, আর !!{elimination} বিধিরগুলি
“বিনাশক”। কোনো বিধির যতগুলি পূর্বধারণা,
সংশ্লিষ্ট অপেক্ষকচিহ্ন ততগুলি যুক্তি নেয়।
যেমন, \Intro{\land} ও \Elim{\land}[1]-এর
অপেক্ষকচিহ্ন যথাক্রমে $c^\land$
(দ্বিস্থানী) এবং $d^\land_1$ (একস্থানী)
হতে পারে। \Log{tN2i}-তে বিধিদুটি তখন
\[
\Axiom$\Gamma_1 \fCenter N: !A$
\Axiom$\Gamma_2 \fCenter M: !B$
\RightLabel{\Intro{\land}}
\BinaryInf$\Gamma_1, \Gamma_2 \fCenter c^\land(N, M) : !A \land !B $
\DisplayProof
\quad
\Axiom$\Gamma \fCenter N: !A \land !B$
\RightLabel{\Elim{\land}[1]}
\UnaryInf$\Gamma \fCenter d^\land_1(N) : !A$
\DisplayProof
\]
% BN-SRC-612 TARGET-CORRECTION-BEGIN
\par\smallskip\noindent\textnormal{\small\textbf{উৎসসংশোধন (BN-SRC-612):} দ্বিতীয় পূর্বধারণার পদ \(M\); সংযোজন-নির্মাতার ফলেও উৎসের ছোট হাতের \(m\) বদলে \(M\) রাখা হয়েছে।}
% BN-SRC-612 TARGET-CORRECTION-END
কোনো অনুমানবিধি একটি অনুমান অবমুক্ত করলে,
অর্থাৎ বিধিটির কোনো পূর্বধারণায় চিহ্নিত
!!{formula}~$x:!A$ থাকে অথচ সিদ্ধান্তের
প্রসঙ্গে সেটি থাকে না, প্রমাণপদেও তা
দেখাতে হবে। এমন বিধির সংশ্লিষ্ট
অপেক্ষকচিহ্নগুলি সংশ্লিষ্ট চলরাশি
\emph{বদ্ধ করে}। যেমন, \Intro{\lif}
বিধি $x:!A, \Gamma \Sequent !B$
পূর্বধারণা থেকে $\Gamma \Sequent !A \lif !B$
দেয়। নির্মাতা $c^{\lif}$ একটি যুক্তি
নেয়, কিন্তু চলরাশি~$x$ বদ্ধ হয়েছে
তাও দেখাতে হয়। এভাবেই প্রমাণপদ
প্রকাশ করে যে $x$-চিহ্নিত অনুমানটি
অবমুক্ত হয়েছে। যেমন, পূর্বধারণার
প্রমাণপদ~$N$ হলে সিদ্ধান্তের প্রমাণপদ
$c^{\lif}([x]N)$ লিখতে পারি; বিধিটি
\[
\Axiom$x:!A, \Gamma \fCenter N: !B$
\RightLabel{\Intro{\lif}}
\UnaryInf$\Gamma \fCenter c^{\lif}([x]N) : !A \lif !B$
\DisplayProof
\]

!!a{proof} সম্পূর্ণ পুনর্গঠন করতে
আরও কিছু তথ্য চাই। কোনো বিধির
পূর্বধারণার !!{formula}s থেকে সিদ্ধান্তের
!!{formula} পুরোপুরি নির্ধারিত না হলে
সেই তথ্য পদে যোগ করতে হয়।
\Intro{\lif}-এ এমন হয়; তাই
$c^{\lif}([x:!A]N)$ লিখব। \Intro{\lor}
ও~\FalseInt-এও এমন হয়। তাই সংশ্লিষ্ট
তথ্যবাহী অপেক্ষকচিহ্ন ব্যবহার করি; যেমন,
\[
\Axiom$\Gamma \fCenter N:!A$
\RightLabel{\Intro{\lor}[1]}
\UnaryInf$\Gamma \fCenter c^{\lor{!B}}_1(N):!A \lor !B$
\DisplayProof
\quad
\Axiom$\Gamma \fCenter N:\lfalse$
\RightLabel{\FalseInt}
\UnaryInf$\Gamma \fCenter d^\lfalse_{!A}(N):!A$
\DisplayProof
\]
% BN-SRC-613 TARGET-CORRECTION-BEGIN
\par\smallskip\noindent\textnormal{\small\textbf{উৎসসংশোধন (BN-SRC-613):} \(\lor\)-ভূমিকা-বিধির এক পূর্বধারণার প্রমাণপদ \(N\) উৎসের সিদ্ধান্তে বাদ ছিল; নির্মাতার যুক্তি হিসেবে \((N)\) যোগ করা হয়েছে।}
% BN-SRC-613 TARGET-CORRECTION-END

এই ধারণাগুলি দিয়ে সিকোয়েন্ট-রীতির
এমন একটি স্বাভাবিক নিষ্পাদন ব্যবস্থা
গড়া যায়, যেখানে প্রত্যেক সিকোয়েন্টের
রূপ $\Gamma \Sequent N : !A$ এবং $N$
একটি \emph{প্রমাণপদ}।

এমন প্রমাণপদ ও তথাকথিত
\emph{টাইপযুক্ত ল্যাম্বডা ক্যালকুলাস}-এর
পদের মধ্যে ঘনিষ্ঠ সম্পর্ক আছে। সেই সম্পর্ক
মেনে উপরের সাধারণ নির্মাতা-বিনাশক চিহ্নের
বদলে ল্যাম্বডা ক্যালকুলাসের প্রচলিত
চিহ্ন ব্যবহার করব। যেমন,
$c^{\lif}([x:!A]N)$-এর বদলে
$\lambd[\typeof{x}{!A}][N]$,
$d^{\lif}(N,M)$-এর বদলে $(NM)$,
$c^\land(N,M)$-এর বদলে $\pair{N}{M}$,
আর $d^\land_i(N)$-এর বদলে
$\proj{i}{N}$ লিখব।

\begin{defn}
  নিচের বিধিগুলি দিয়ে আরোহক্রমে
  প্রমাণপদ গঠিত হয়:
  \begin{enumerate}
  \item প্রত্যেক চলরাশি $x$ একটি
    প্রমাণপদ (স্বতঃসিদ্ধ)।
  \item $x$ চলরাশি, $!A$ !!a{formula}
    এবং $N$ প্রমাণপদ হলে
    $\lambd[\typeof{x}{!A}][N]$-ও
    প্রমাণপদ~(\Intro\lif)।
  \item $N$ ও $M$ প্রমাণপদ হলে $(NM)$-ও
    প্রমাণপদ~(\Elim\lif)।
  \item $N$ ও $M$ প্রমাণপদ হলে
    $\pair{N}{M}$ প্রমাণপদ~(\Intro\land)।
  \item $N$ প্রমাণপদ হলে
    $\proj{i}{N}$-ও প্রমাণপদ; এখানে
    $i$ হলো $1$ অথবা~$2$~(\Elim\land[i])।
  \item $N$ প্রমাণপদ এবং $!A$ সূত্র হলে
    $\inj[!A]{i}{N}$ প্রমাণপদ; এখানে
    $i$ হলো $1$ অথবা~$2$~(\Intro\lor[i])।
% BN-SRC-614 TARGET-CORRECTION-BEGIN
\par\smallskip\noindent\textnormal{\small\textbf{উৎসসংশোধন (BN-SRC-614):} এই ধারায় \(N\) প্রমাণপদ হিসেবে বাঁধা; উৎসের \(\inj[!A]{i}{!M}\)-এ অনির্ধারিত \(M\) ও সূত্রচিহ্নের অতিরিক্ত বিস্ময়চিহ্ন ছিল। অনুরূপ tN2i বিধিতে \(\inj[!B]{1}{N}\) আছে।}
% BN-SRC-614 TARGET-CORRECTION-END
  \item $M$, $N_1$, $N_2$ প্রমাণপদ,
    আর $x_1$, $x_2$ চলরাশি হলে
    $\dcase{M}{x_1}{N_1}{x_2}{N_2}$
    প্রমাণপদ~(\Elim\lor)।
  \item $N$ প্রমাণপদ এবং $!A$
    !!a{formula} হলে
    $\abort{!A}{N}$ প্রমাণপদ~(\FalseInt)।
  \end{enumerate}
\end{defn}

প্রত্যেক প্রমাণপদ কোনো !!a{proof}-এর
প্রমাণপদ নয়। যেমন, $\pair{N}{M}$ পদ
\Intro{\land} বিধিতে শেষ হওয়া প্রমাণকে
সংকেতায়িত করে; তাই তার অন্তিম
!!{formula} হলো সংযোজন~$!A \land !B$।
এটি \Elim{\lif} বিধির প্রধান পূর্বধারণা
হতে পারে না; ফলে
$(\pair{N}{M}P)$ কোনো সঠিক প্রমাণের
প্রমাণপদ হতে পারে না। শুধু
\Log{tN2ip}-এর বিধি মেনে চলা
প্রমাণপদগুলিই \emph{সঠিক} প্রমাণপদ।
বিধিগুলি \olref{tab:tN2ip}-এ দেওয়া আছে।

\subfile{rules-tN2}

\end{document}
